\section*{Chapter \ref{ch:g7:strings} --- Strings}

\begin{solution}{exo:g7:strings:1}
\code{s[1:4]} is \code{"omp"} (indices 1, 2, 3), \code{s[-1]} is \code{"r"} and
\code{s[::-1]} is \code{"retupmoc"}.
\end{solution}

\begin{solution}{exo:g7:strings:2}
The \omterm{def:g7:lists:list}{list} \code{["a", "b", "c"]}, and the \omterm{def:g7:strings:string}{string} \code{"1 + 2 + 3"}.
\end{solution}

\begin{solution}{exo:g7:strings:3}
The capitals follow \code{A} (65) in order, so \code{ord("C")} is 67. The small
letters follow \code{a} (97): 100 is \code{"d"}.
\end{solution}

\begin{solution}{exo:g7:strings:4}
\omcode{code/grade-7/03-strings/exercises/spaces.py}{1}{6}{Counting
spaces.}\label{lst:g7:strings:sol-spaces}

It prints \code{6 6}: both ways agree.
\end{solution}

\begin{solution}{exo:g7:strings:5}
The indices of a \omterm{def:g7:strings:string}{string} of length $n$ go from 0 to $n - 1$, so \omterm{def:g7:lists:list}{index}
\code{len(s)} is past the end. The last \omterm{def:g7:strings:string}{character} is \code{s[len(s) - 1]}, or
simply \code{s[-1]}.
\end{solution}

\begin{solution}{exo:g7:strings:6}
\code{s.lower()} is compared with \code{s[::-1]}, the reverse of the original
\omterm{def:g7:strings:string}{string}, which still has its capital: \code{"kayak"} against \code{"kayaK"}. Both
sides must be in small letters; the simplest way is to change \code{s} first:

\omcode{code/grade-7/03-strings/exercises/palindrome_fixed.py}{1}{6}{Small
letters on both sides: it prints \code{True}.}\label{lst:g7:strings:sol-palindrome}
\end{solution}

\begin{solution}{exo:g7:strings:7}
\omcode{code/grade-7/03-strings/exercises/initials.py}{1}{9}{Initials.}\label{lst:g7:strings:sol-initials}

It prints \code{A.L.} and \code{G.B.H.}.
\end{solution}

\begin{solution}{exo:g7:strings:8}
\omcode{code/grade-7/03-strings/exercises/star_vowels.py}{1}{8}{Hiding the
vowels.}\label{lst:g7:strings:sol-star}

It prints \code{*nc* *p*n * t*m*}. (\code{text.replace} would need one call
per vowel, capitals included.)
\end{solution}

\begin{solution}{exo:g7:strings:9}
The \omterm{def:g7:strings:slice}{slice} starts at \omterm{def:g7:lists:list}{index} $i$ and ends before $i + 3$, so its last \omterm{def:g7:strings:string}{character}
has \omterm{def:g7:lists:list}{index} $i + 2$, at most 9: $i$ goes from 0 to 7, which makes 8 \omterm{def:g7:strings:slice}{slices}.
\end{solution}

\begin{solution}{exo:g7:strings:10}
A best-so-far walk, comparing lengths:

\omcode{code/grade-7/03-strings/exercises/longest_word.py}{1}{10}{The longest
word.}\label{lst:g7:strings:sol-longest}

It prints \code{quick}: \code{brown} and \code{jumps} are as long, but come
later, and only a strictly longer word replaces the best.
\end{solution}

\begin{solution}{exo:g7:strings:11}
\code{s[::-1]} puts the last \omterm{def:g7:strings:string}{character} first, and so on; reversing again puts
every \omterm{def:g7:strings:string}{character} back in its place. \code{s[:k]} holds the \omterm{def:g7:strings:string}{characters} with
indices 0 to $k - 1$, and \code{s[k:]} those with indices $k$ to the end: each
\omterm{def:g7:lists:list}{index} is in exactly one of the two \omterm{def:g7:strings:slice}{slices}, in order, so gluing them gives
\code{s}. With $k = 0$ the first \omterm{def:g7:strings:slice}{slice} is empty, and with $k = \code{len(s)}$
the second one is.
\end{solution}

\begin{solution}{exo:g7:strings:12}
\omcode{code/grade-7/03-strings/exercises/mirror.py}{1}{13}{A mirror
test.}\label{lst:g7:strings:sol-mirror}

\code{bid}, \code{bud} and \code{pq} look the same in the mirror; \code{pod}
becomes \code{boq}, and \code{box} becomes \code{xod}.
\end{solution}

\begin{solution}{pb:g7:strings:1}
\textbf{1.} \code{02022020} and \code{12022021}: reading them from the right
gives the same digits.

\textbf{2.} The year $2\,0\,c\,d$ reversed is $d\,c\,0\,2$: the month is always
\code{02}, February, and the day is $d\,c$, the last two digits of the year in
reverse order.

\textbf{3.} For 2092 the day is 29: the 29th of February 2092. It exists, since
2092 is divisible by 4 and is a leap year.

\textbf{4.} \code{"0" + str(n)} puts a 0 in front: \code{"07"} or \code{"012"}.
The \omterm{def:g7:strings:slice}{slice} \code{[-2:]} keeps the last two \omterm{def:g7:strings:string}{characters}: \code{"07"} and
\code{"12"}.

\textbf{5.} The year already has four digits, for every year from 1000 to
9999.

\textbf{6.} The leap-year rule of \cref{ch:g6:if-else}: February has 29 days in
a leap year and 28 otherwise, and all the palindrome dates of this century are
in February.

\textbf{7.} It agrees:

\omoutput{code/grade-7/03-strings/outputs/dates.txt}{The two dates of
question 1.}\label{out:g7:strings:dates}

\textbf{8.} Three \omterm{def:g6:for-loops-turtle:nested}{nested loops} (year, month, day), each date turned into a
\omterm{def:g7:strings:string}{string} and tested:

\omcode{code/grade-7/03-strings/listings/all_dates.py}{1}{13}{Every palindrome
date of the century.}\label{lst:g7:strings:all-dates}

\textbf{9.} 29 dates, from 10 February 2001 to 29 February 2092:

\omoutput{code/grade-7/03-strings/outputs/all_dates.txt}{The palindrome dates
from 2001 to 2099.}\label{out:g7:strings:all-dates}

\textbf{10.} A year $2\,0\,c\,d$ has one palindrome date when the day $d\,c$
exists in February, and none otherwise: there is only one possible month and
day for each year. The day $d\,c$ exists when $d$ is 0 (and $c$ is not 0), when
$d$ is 1, or when $d$ is 2 and $c$ is at most 8, or 9 in a leap year. So the
years ending in 0 (except 2000), in 1, and in 2 up to 2082, plus 2092: $9 + 10 +
10 = 29$ years.

\textbf{11.} It is the only one that falls on the 29th of February, a day that
exists only in leap years.

\textbf{12.} \textbf{Result.} There are 29 palindrome dates between 2001 and
2099, all in February, from 10-02-2001 to 29-02-2092.
\end{solution}
